\section{Rodando o primeiro container e comandos de sobrevivência}

Após a instalação, é possível rodar o primeiro container com o seguinte comando:

\begin{verbatim}
docker run -d -p 8080:80 docker/welcome-to-docker
\end{verbatim}

O \textit{frontend} dessa aplicação fica acessível na porta 8080, pelo endereço \texttt{http://localhost:8080} no navegador. Após a execução, o container aparece como rodando e a imagem passa a constar nas images locais.

\subsection{Comandos básicos}

\subsubsection{Executando containers}

\begin{verbatim}
docker run -d -p 8080:80 docker/welcome-to-docker
\end{verbatim}

A flag \texttt{-d} executa o container em modo \textit{detached} (em segundo plano). A flag \texttt{-p 8080:80} mapeia a porta 8080 do host para a porta 80 do container.

\subsubsection{Listando containers}

\begin{verbatim}
docker ps          # lista containers em execução
docker ps -a       # lista todos (incluindo os parados)
\end{verbatim}

\subsubsection{Parando containers}

\begin{verbatim}
docker ps              # obter o ID do container
docker stop <ID>       # parar o container
\end{verbatim}

Não é necessário digitar o ID completo --- basta os caracteres suficientes para torná-lo único entre todos os containers existentes (inclusive os parados).

